% SPDX-License-Identifier: Apache-2.0
\section{Reading Memory in Bursts}
\label{sec:x-dma}

A peripheral that wants a lot of memory does not want the core to
fetch it. The video peripheral holds its framebuffer inside the unit
and the core writes it a pixel at a time over AXI-Lite: 19\,200
register writes for a frame of the board's 160 by 120, 307\,200 for
a full 640 by 480, and a processor whose
whole job during them is copying. The Ethernet peripheral is the same
shape a byte at a time.

The alternative is for the peripheral to be a host on the bus. The
core says where the data is and how much of it there is; the
peripheral fetches it. \code{LineFetch} in \code{//lib/parts} is that,
and it is the half the video and the MAC both want: given an address
and a count it issues read bursts and puts the words that come back
into a channel. What consumes the channel is the peripheral's
business. For the video it crosses to the pixel clock through the
part of \S\ref{sec:x-cdc} and feeds the raster.

\lstinputlisting[rangebeginprefix=//\ begin\{,rangebeginsuffix=\},%
rangeendprefix=//\ end\{,rangeendsuffix=\},includerangemarker=false,%
linerange=state-state,caption={\code{lib/parts/src/dma.rs}: what it
remembers between cycles.},label={lst:x-dma}]{lib/parts/src/dma.rs}

\subsection{Two things the bus does backwards}

Both of these were written the other way first, and both deadlock.

\emph{The identifier arrives with the burst, not before it.} A host
takes an \code{Issue}, chooses a free identifier and says which one on
\code{grant}. So a client sends the burst and catches the identifier
as it comes back. Asking for an identifier and waiting for it is a
deadlock: nothing is granted until something is issued.

\emph{A read burst is finished by its last beat, not by \code{done}.}
The host puts a write's response on \code{done} and answers a read on
the beat channel. So the identifier goes back as the last beat is
taken, and the last beat is only taken when there is room to give it
back, which is why the room for the release is part of the condition
for accepting the beat.

One burst is in flight at a time, deliberately. A second would want a
second identifier, a tracker to tell the two apart, and an answer for
what happens when they come back out of order. None of that buys
anything while the consumer is a raster taking one pixel per visible
cycle.

\subsection{The gain, as something a test can fail}

The run fetches sixty-four words through sixteen-beat bursts, across a
real link with \code{AxiHost}, \code{AxiPer} and a memory on the far
side. Four address phases carry what would otherwise take sixty-four.

It asserts rather than prints. Every word arrived. Each word is the
one at its own address, because the memory holds its address in each
word, so a word from the wrong burst or out of order says so. And the
run took fewer cycles than a path spending an address phase per word
could manage: sixty-four words in 93 cycles, against a bound of 128.

That last one is the point of the example. A throughput number in a
log is a number somebody has to read and believe; a bound a run must
come in under is a claim the build checks every time.

The memory model in the example answers bursts rather than words, and
that is not incidental. A burst arrives as one request carrying
\code{len}, beats less one, and is owed \code{len + 1} beats with the
last of them marked. A model that answered one beat per request would
make a burst look like a word: the words would still all arrive, the
cycle count would still pass, and the example would demonstrate
nothing about bursting. Getting that wrong is the way this example
could have been green and worthless.

\lstinputlisting[style=ascii,caption={What \code{ex\_dma} printed when
the build ran it.},label={lst:x-dma-out}]{out_dma.txt}
